\documentclass[11pt,a4paper]{article}

% ---------------------------------------------------------------------------
% Paquetes
% ---------------------------------------------------------------------------
\usepackage[utf8]{inputenc}
\usepackage[T1]{fontenc}
\usepackage[spanish,es-tabla]{babel}
\usepackage[a4paper,margin=2.5cm]{geometry}
\usepackage{graphicx}
\usepackage{booktabs}
\usepackage{array}
\usepackage{longtable}
\usepackage{caption}
\usepackage{xcolor}
\usepackage{enumitem}
\usepackage{amsmath}
\usepackage{hyperref}
\usepackage{fancyhdr}

\hypersetup{
    colorlinks=true,
    linkcolor=blue!50!black,
    urlcolor=blue!50!black,
    citecolor=blue!50!black,
    pdftitle={Detección de tanques de almacenamiento de petróleo: métodos deterministas de OpenCV frente a una CNN},
    pdfauthor={},
}

\graphicspath{{results/figures/}}

\setlength{\headheight}{14pt}
\pagestyle{fancy}
\fancyhf{}
\fancyhead[L]{Fundamentos de la Visión por Computador}
\fancyhead[R]{\thepage}
\renewcommand{\headrulewidth}{0.4pt}

\setlength{\parskip}{0.5em}
\setlength{\parindent}{0pt}

% ---------------------------------------------------------------------------
% Datos del documento
% ---------------------------------------------------------------------------
\title{
    \textbf{Fundamentos de la Visión por Computador} \\[0.5em]
    \large Detección de tanques de almacenamiento de petróleo en imágenes satelitales:\\
    métodos deterministas de OpenCV frente a una CNN (YOLOv8)
}
\author{Gotzon Viteri Fuente}
\date{\today}

% ---------------------------------------------------------------------------
\begin{document}

\maketitle
\thispagestyle{fancy}

\tableofcontents
\newpage

% =============================================================================
\section{Introducción y Objetivo General}
% =============================================================================

\subsection{Definición del problema}

El problema abordado es la \textbf{detección de tanques de almacenamiento de
petróleo} en imágenes de satélite de alta resolución. Dada una imagen aérea de una
instalación industrial o portuaria, el sistema debe localizar cada tanque presente
mediante una caja delimitadora (\emph{bounding box}), sin necesidad de clasificar
subtipos de tanque ni estimar su volumen.

Los tanques de almacenamiento son estructuras cilíndricas que, vistas desde arriba,
aparecen como \textbf{círculos o elipses casi perfectos}, de tamaño relativamente
uniforme dentro de una misma instalación pero muy variable entre instalaciones
distintas (de pocos metros a más de 100 m de diámetro). Se agrupan característicamente
en ``granjas de tanques'' --- filas y rejillas muy densas de decenas o cientos de
unidades --- lo que convierte la \textbf{separación de instancias contiguas} en el
principal reto técnico del problema, por encima de la propia detección de la forma
circular.

Este tipo de detección tiene aplicaciones reales en inteligencia geoespacial:
estimación de reservas estratégicas de petróleo a partir de la sombra proyectada por
el techo flotante de los tanques, monitorización de actividad industrial, o
vigilancia de infraestructuras críticas --- todas ellas tareas que hoy dependen de
imagen satelital comercial y que motivaron a Airbus Defense and Space a publicar el
dataset utilizado en este proyecto.

\subsection{Estado del arte y trabajos relacionados}

Se han identificado y consultado las siguientes referencias directamente relacionadas
con este problema y, en varios casos, con este mismo dataset:

\begin{enumerate}
    \item \textbf{Hough Circle Transform} \cite{duda1972hough} --- la técnica clásica
    de referencia para detectar formas circulares mediante acumulación de votos en un
    espacio de parámetros (centro, radio), implementada en OpenCV como
    \texttt{cv2.HoughCircles}. Es el punto de partida natural para cualquier detector
    de objetos circulares, pero \textbf{no discrimina textura de fondo}: en escenas
    con muchos bordes (carreteras, parcelas agrícolas) genera un número inmanejable de
    falsos círculos si se ajusta para ser sensible a objetos pequeños.

    \item \textbf{Faudi, J.} --- \emph{``Oil Storage Detection on Airbus Imagery with
    YOLOX''} \cite{faudi_yolox}, curador del propio dataset. Trocea las imágenes de
    2560$\times$2560 en \emph{tiles} de 640$\times$640 con solape de 64 px (protocolo
    que hemos replicado en este proyecto) y entrena YOLOX-s durante solo 10 épocas,
    obteniendo AP@0.5 = 0.856. \emph{Limitación}: no compara contra ningún método
    clásico ni analiza en qué tipo de escena falla el modelo.

    \item \textbf{Ouzounis, G.} --- \emph{``Oil-Storage Tank Instance Segmentation
    with Mask R-CNN''} \cite{ouzounis_maskrcnn}. Convierte las cajas del mismo dataset
    en máscaras circulares sintéticas y entrena Mask R-CNN sobre recortes de
    128$\times$128. Reporta cualitativamente el mismo problema estructural que hemos
    encontrado nosotros: ``leakage'' (fuga) al separar tanques contiguos por falta de
    contexto en recortes pequeños. \emph{Limitación}: no aporta métricas cuantitativas
    de detección, solo observaciones cualitativas.

    \item \textbf{Wang et al.} --- \emph{``Storage tank detection in remote sensing
    images based on circular bounding boxes and large selective kernel''}
    \cite{wang2025circular}, Scientific Reports (2025). Propone sustituir la caja
    rectangular estándar de un detector tipo YOLO por una parametrización circular
    explícita (centro + radio) y un kernel selectivo de campo receptivo amplio,
    alineando mejor la representación del modelo con la forma real del objeto. Es la
    línea de mejora más prometedora que hemos identificado para superar el techo de
    F1$\approx$0.84 obtenido con una caja rectangular genérica, aunque no hemos podido
    confirmar si sus experimentos usan este mismo dataset (artículo con acceso
    restringido).

    \item \textbf{Beucher, S. \& Meyer, F.} --- algoritmo de segmentación
    \emph{watershed} controlado por marcadores \cite{beucher1993watershed} (técnica
    clásica de procesamiento morfológico, disponible en \texttt{cv2.watershed}), el
    método de referencia en visión por computador clásica para separar objetos
    convexos que se tocan (su aplicación más citada es la separación de células en
    microscopía). Es la base de uno de los experimentos deterministas de este
    proyecto.
\end{enumerate}

\textbf{Ningún trabajo consultado reporta un \emph{baseline} puramente determinista}
(sin aprendizaje) para este dataset con el que comparar cuantitativamente. Esa
ausencia es precisamente el hueco que este proyecto cubre.

\subsection{Valor añadido del proyecto}

Frente a los antecedentes anteriores, este proyecto no se limita a entrenar un
detector más: aporta

\begin{itemize}
    \item Un \textbf{barrido sistemático de 7 técnicas deterministas} de complejidad
    creciente (Hough circular, Hough tileado, contornos + circularidad, \emph{blob
    detection}, segmentación \emph{watershed} controlada por marcadores, detección
    multiescala Diferencia-de-Gaussianas) más un post-proceso de filtrado por
    densidad --- todas evaluadas bajo el \textbf{mismo protocolo exacto} (mismo
    \emph{split}, mismo umbral de IoU, mismo código de métricas) que el modelo de
    referencia basado en CNN, algo que ninguna de las referencias anteriores ofrece.

    \item Una \textbf{cuantificación del techo real del enfoque clásico} en este
    problema concreto (F1$\approx$0.354, frente a F1$\approx$0.843 de una CNN
    \emph{fine-tuned}) y, más importante, un \textbf{diagnóstico de \emph{por qué}}
    se produce ese techo: no es un problema de ajuste de hiperparámetros (se han
    probado más de 150 configuraciones distintas en total) sino una limitación
    estructural --- los detectores de forma clásicos razonan objeto a objeto a partir
    de bordes o umbrales locales, y el 50\,\% de las imágenes del dataset contienen
    más de 100 tanques colocados en rejillas muy densas, el escenario exacto en el
    que separar instancias contiguas sin contexto de alto nivel es más difícil.

    \item La confirmación (Sección~\ref{sec:validacion}) de que nuestro resultado con
    YOLOv8 es coherente con el resultado publicado por el propio curador del dataset
    usando YOLOX y un protocolo de \emph{tiling} casi idéntico, lo que valida la
    metodología de evaluación seguida.
\end{itemize}

% =============================================================================
\section{Metodología y Desarrollo}
% =============================================================================

\subsection{Conjunto de datos}

Se ha utilizado el dataset \textbf{Airbus Oil Storage Detection}
\cite{airbus_dataset} (Airbus DS Intelligence / Scale AI), publicado bajo licencia
Creative Commons BY-NC-SA 4.0.

\begin{table}[h!]
\centering
\begin{tabular}{@{}ll@{}}
\toprule
\textbf{Propiedad} & \textbf{Valor} \\
\midrule
Nº de imágenes & 98 (+ 5 adicionales sin anotar en \texttt{extras/}) \\
Resolución de imagen & 2560 $\times$ 2560 px \\
Resolución de suelo & $\sim$1.2--1.5 m/píxel (satélite SPOT) \\
Condiciones de adquisición & Óptico, luz natural diurna, ubicaciones \\
 & mundiales diversas (iluminación/sombra variable), \\
 & imagen fija ortorectificada (sin movimiento) \\
Nº total de anotaciones & 13\,592 cajas, clase única \texttt{oil-storage-tank} \\
Tanques por imagen & mín. 14, mediana 99.5, máx. 893 \\
Tamaño de tanque (lado de caja) & mín. 3 px, mediana 18 px, máx. 108 px \\
\bottomrule
\end{tabular}
\caption{Características del dataset Airbus Oil Storage Detection.}
\label{tab:dataset}
\end{table}

\textbf{Limpieza de datos}: se detectó que algunas anotaciones del CSV tienen
$y_2 < y_1$ (altura negativa, p.\ ej.\ una fila con $h = -77$), por lo que el
\emph{parser} de anotaciones normaliza cada caja tomando
$(\min(x_1,x_2), \min(y_1,y_2), \max(x_1,x_2), \max(y_1,y_2))$ en vez de asumir que
las coordenadas ya vienen ordenadas.

\textbf{Partición}: \emph{split} fijo del 80/20 a nivel de imagen completa (no de
anotación individual, para evitar fuga de información entre \emph{train} y
\emph{val}), con semilla aleatoria fija (42), generando 78 imágenes de
entrenamiento/ajuste y 20 de validación. El mismo \emph{split} se reutiliza para
\textbf{todos} los métodos comparados, de forma que los números de la
Sección~\ref{sec:resultados} son directamente comparables entre sí.

\subsection{Arquitectura del proyecto}

El proyecto se organiza en módulos reutilizables (\texttt{src/}) y \emph{scripts}
ejecutables (\texttt{scripts/}):

\begin{itemize}[leftmargin=2.2cm]
    \item[\texttt{src/dataset.py}] carga/parseo de anotaciones, \emph{split}
    train/val, lectura/escritura de imagen segura en Windows.
    \item[\texttt{src/evaluate.py}] IoU, \emph{matching} \emph{greedy}, métricas
    P/R/F1, NMS genérico.
    \item[\texttt{src/tiling.py}] partición de una imagen grande en \emph{tiles} con
    solape.
    \item[\texttt{src/hough\_detector.py}] Hough Circles sobre imagen completa +
    \emph{grid search}.
    \item[\texttt{src/classic\_methods.py}] 7 variantes deterministas adicionales,
    todas tileadas.
    \item[\texttt{src/yolo\_data.py}] conversión del dataset a formato YOLO (tiling +
    etiquetas).
    \item[\texttt{src/yolo\_detector.py}] entrenamiento e inferencia tileada de
    YOLOv8.
    \item[\texttt{src/visualize.py}] superposición de cajas GT/predicción sobre la
    imagen.
\end{itemize}

Los \emph{scripts} de \texttt{scripts/} orquestan el ajuste de hiperparámetros
(\emph{grid search}), la evaluación sobre validación y la generación de tablas y
figuras comparativas para cada método.

\textbf{Librerías utilizadas}: OpenCV (\texttt{opencv-python}) para todo el
procesamiento de imagen de bajo/medio nivel (filtrado, umbralización, transformada de
Hough, \texttt{SimpleBlobDetector}, morfología, \emph{watershed}), NumPy para álgebra
vectorizada (IoU, NMS), pandas para el parseo de anotaciones, \texttt{scikit-image}
\cite{vanderwalt2014scikit} para detección de \emph{blobs} multiescala y
\texttt{peak\_local\_max}/\texttt{regionprops}, SciPy para el etiquetado de
componentes conexas, \texttt{ultralytics} (YOLOv8 \cite{yolov8}, sobre PyTorch con
aceleración CUDA) como referencia de red neuronal, y matplotlib/OpenCV
\cite{bradski2000opencv} para las visualizaciones.

\subsection{Protocolo de evaluación común}

Para que la comparación entre técnicas tan distintas sea honesta, todas comparten:

\begin{itemize}
    \item \textbf{Emparejamiento \emph{greedy} por IoU}: cada predicción se empareja,
    por orden de confianza (o de llegada si el detector no aporta confianza, como los
    métodos clásicos), con la anotación real no usada de mayor solapamiento,
    aceptando el emparejamiento solo si $\text{IoU} \geq 0.5$.
    \item \textbf{Métricas}: precisión, \emph{recall}, F1 (métrica principal para
    elegir hiperparámetros), IoU medio de los aciertos (mide \emph{calidad de
    localización} independientemente de cuántos aciertos haya) y tiempo de
    inferencia por imagen.
    \item \textbf{Ajuste de hiperparámetros exclusivamente sobre \emph{train}}: cada
    método se afina por \emph{grid search} maximizando F1 sobre una muestra de
    entrenamiento (8--12 imágenes según el experimento) y la configuración ganadora
    se congela antes de evaluarse, una única vez, sobre las 20 imágenes de
    validación.
    \item \textbf{NMS genérico} para fusionar detecciones solapadas entre
    \emph{tiles} contiguos.
\end{itemize}

\subsection{Técnicas deterministas evaluadas}

Todas operan por \textbf{teselado (\emph{tiling})} de la imagen en ventanas de
640$\times$640 px con 64 px de solape (salvo el primer experimento, que se hace sobre
la imagen completa como punto de partida), replicando la misma estrategia de
partición usada para YOLO, de forma que la comparación no esté sesgada por trabajar a
distinta escala espacial.

\begin{enumerate}
    \item \textbf{Hough Circles, imagen completa} --- \texttt{cv2.HoughCircles} sobre
    la imagen de 2560$\times$2560 tras realce con CLAHE y filtro de mediana.
    \emph{Grid search} sobre \texttt{dp}, \texttt{param1}, \texttt{param2},
    \texttt{minDist}.
    \item \textbf{Hough Circles, tileado} --- la misma transformada, pero corrida
    ventana a ventana y fusionando resultados con NMS.
    \item \textbf{Contornos + circularidad + \emph{watershed} local} --- umbral
    adaptativo, contornos filtrados por circularidad
    ($4\pi\,\text{Área}/\text{Perímetro}^2$), con separación por \emph{watershed}
    solo cuando un contorno es sospechosamente grande (probable fusión de varios
    tanques).
    \item \textbf{\texttt{cv2.SimpleBlobDetector}} --- detector de \emph{blobs}
    nativo de OpenCV, configurado para buscar regiones estables a través de
    múltiples umbrales de brillo, filtrando por área, circularidad, convexidad y
    polaridad de color (\emph{blobs} claros vs.\ oscuros).
    \item \textbf{Segmentación \emph{watershed} universal controlada por
    marcadores} --- umbral global por percentil de brillo, \emph{distance
    transform}, máximos locales como semillas de primer plano, fondo seguro por
    dilatación, y \texttt{cv2.watershed} sobre toda la máscara (no solo sobre
    \emph{blobs} sospechosos, a diferencia de la técnica 3).
    \item \textbf{Detección multiescala Diferencia-de-Gaussianas (DoG)} ---
    \texttt{skimage.feature.blob\_dog} \cite{lindeberg1998scale}, que barre un rango
    continuo de escalas ($\sigma$) en vez de trabajar a escala fija.
    \item \textbf{Post-proceso: filtro por densidad de vecinos} --- aplicado sobre
    el mejor detector (técnica 4), descarta detecciones sin ningún otro candidato
    cercano, bajo la hipótesis de que los tanques reales aparecen casi siempre en
    grupo.
\end{enumerate}

El detalle experimento a experimento --- incluyendo un bug real de diseño detectado
y corregido durante el desarrollo (Sección~\ref{sec:bug}) --- está documentado
íntegramente en el fichero \texttt{EXPERIMENTS.md} que acompaña al código fuente.

\subsection{Modelo de referencia: YOLOv8}

Como referencia de lo que consigue un enfoque de aprendizaje profundo con el mismo
volumen de datos, se entrenó \textbf{YOLOv8n} (variante más pequeña de Ultralytics,
preentrenada en COCO) mediante \emph{fine-tuning}:

\begin{itemize}
    \item Las imágenes de entrenamiento se trocean en \emph{tiles} de
    640$\times$640 px con 64 px de solape (mismo esquema que los métodos
    clásicos), conservando solo los \emph{tiles} con al menos un tanque, y
    descartando cajas con menos del 40\,\% de su área dentro del \emph{tile}.
    \item Entrenamiento: 60 épocas, tamaño de imagen 640, sobre GPU NVIDIA RTX
    3060 (CUDA).
    \item Inferencia sobre las imágenes de validación: se aplica el mismo
    teselado, se corre el modelo \emph{tile} a \emph{tile} y se fusionan las
    detecciones con NMS, reproyectando a las coordenadas de la imagen original de
    2560$\times$2560 --- el mismo procedimiento, con el mismo código de teselado,
    que usan los métodos clásicos.
\end{itemize}

% =============================================================================
\section{Resultados y Conclusiones}
\label{sec:resultados}
% =============================================================================

\subsection{Tabla comparativa (20 imágenes de validación)}

\begin{table}[h!]
\centering
\small
\begin{tabular}{@{}lccccc@{}}
\toprule
\textbf{Método} & \textbf{Precisión} & \textbf{Recall} & \textbf{F1} & \textbf{IoU medio} & \textbf{Tiempo/img} \\
\midrule
Hough Circles, imagen completa & 0.113 & 0.044 & 0.063 & 0.814 & 0.073 s \\
DoG multiescala & 0.031 & 0.052 & 0.039 & 0.565 & 3.363 s \\
Contornos + circularidad + \emph{watershed} local & 0.144 & 0.107 & 0.123 & 0.743 & 0.193 s \\
Hough Circles, tileado & 0.443 & 0.104 & 0.169 & 0.769 & 0.175 s \\
\emph{Watershed} universal & 0.287 & 0.217 & 0.247 & 0.657 & 0.794 s \\
\emph{Blob detector} (grid inicial) & 0.323 & 0.291 & 0.306 & 0.671 & 0.439 s \\
\textbf{\emph{Blob detector} (grid refinado)} & \textbf{0.584} & 0.254 & \textbf{0.354} & 0.672 & 0.429 s \\
\emph{Blob detector} + filtro de densidad & 0.626 & 0.248 & 0.355 & 0.671 & 0.428 s \\
\textbf{YOLOv8n \emph{fine-tuned}} & \textbf{0.812} & \textbf{0.876} & \textbf{0.843} & \textbf{0.809} & 0.373 s \\
\bottomrule
\end{tabular}
\caption{Comparativa de todos los métodos evaluados bajo el mismo protocolo, sobre las 20 imágenes de validación.}
\label{tab:resultados}
\end{table}

El mejor método determinista (\emph{blob detector} afinado, con o sin filtro de
densidad) multiplica por \textbf{$\sim$5.6$\times$} el F1 del primer intento con
Hough, pero se queda a \textbf{menos de la mitad} del F1 conseguido con YOLOv8.

\subsection{Análisis de los fallos del sistema determinista}

El hallazgo transversal a los ocho experimentos es que \textbf{el cuello de botella
no es de precisión geométrica sino de \emph{recall}/detección}: el IoU medio de los
aciertos es comparable entre casi todos los métodos (0.56--0.81) e incluso similar al
de YOLO (0.81). Es decir, cuando un método clásico sí encuentra un tanque, dibuja una
caja tan buena como la de la red neuronal. El problema es que \textbf{no encuentra la
mayoría de ellos}: incluso el mejor método clásico solo recupera el 25\,\% de los
tanques reales, frente al 88\,\% de YOLO.

La causa raíz identificada es estructural: los detectores de forma clásicos (Hough,
contornos, \emph{blobs}, \emph{watershed}) razonan \textbf{un objeto a la vez} a
partir de bordes o umbrales locales. Este dataset tiene una mediana de $\sim$100
tanques por imagen colocados en rejillas muy juntas (hasta 893 en la imagen más
densa), y separar instancias que se tocan sin contexto de más alto nivel --- el tipo
de contexto que una CNN aprende directamente de ejemplos anotados --- es precisamente
el escenario donde estas técnicas fallan.

Un segundo patrón, más práctico, es que \textbf{la elección de técnica importa más
que su ajuste fino}: pasar de Hough a \emph{blob detection} (mismo \emph{tiling},
mismo protocolo) duplicó el F1 sin cambiar nada más, simplemente porque el
\texttt{SimpleBlobDetector} incorpora de forma nativa un filtro de \textbf{polaridad
de brillo} (\emph{blobs} claros sobre fondo oscuro) que Hough y DoG no tienen. Los
tanques son, ante todo, los objetos más claros y compactos de la escena; cualquier
técnica que ignore esa señal no supera F1$\approx$0.17.

\subsection{Un fallo de diseño real, encontrado y corregido durante el desarrollo}
\label{sec:bug}

Durante el experimento de \emph{watershed} universal, la primera implementación
devolvía sistemáticamente F1 = 0.000. El diagnóstico reveló que \textbf{nunca se
marcaba explícitamente el fondo} antes de llamar a \texttt{cv2.watershed}: se dejaban
esos píxeles con valor 0 (``desconocido'') en vez de asignarles una etiqueta de fondo
seguro. Sin esa semilla, el algoritmo no tiene desde dónde crecer la región de fondo,
y el 99.97\,\% de las regiones resultantes quedaban reducidas a fragmentos de 1
píxel. La corrección aplicó la receta estándar de \emph{watershed} controlado por
marcadores (fondo seguro por dilatación + frente seguro por los máximos locales +
región desconocida entre ambos), tras lo cual el método pasó a producir resultados
razonables (F1 = 0.247). Se documenta este episodio porque ilustra un error
conceptual no trivial --- y su corrección --- más que un simple problema de
hiperparámetros.

De forma similar, el experimento de DoG resultó ser el peor de todos (F1 = 0.039,
incluso por debajo del Hough más ingenuo) precisamente por \emph{no} incorporar el
filtro de polaridad de brillo que sí tiene el \emph{blob detector} --- confirmando
\emph{post-hoc} la importancia de esa señal.

Por último, el filtro de densidad de vecinos, aplicado como post-proceso sobre el
mejor método, no aportó mejora neta (F1 0.354 $\to$ 0.355): aunque los tanques
reales sí aparecen en grupo, \textbf{los falsos positivos también} (otras
estructuras humanas agrupadas, como naves industriales), así que la densidad
espacial por sí sola no distingue ``cluster de tanques'' de ``cluster de tejados
similares''.

\subsection{Validación externa del resultado con YOLO}
\label{sec:validacion}

El resultado obtenido con YOLOv8n (F1 = 0.843) es consistente con el publicado por
el propio curador del dataset \cite{faudi_yolox}, que con YOLOX-s y un esquema de
\emph{tiling} prácticamente idéntico (640$\times$640, solape de 64 px) obtiene
AP@0.5 = 0.856 tras solo 10 épocas sin ajuste de hiperparámetros. La cercanía de
ambos números, obtenidos de forma completamente independiente, respalda la validez
del protocolo de evaluación seguido en este proyecto.

\subsection{Ejemplo visual}

La Figura~\ref{fig:comparativa} muestra, sobre una misma imagen de validación, las
anotaciones reales (verde) frente a las predicciones de Hough tileado (azul) y de
YOLOv8 (rojo): YOLO cubre casi la totalidad de los tanques reales, mientras que
Hough solo detecta los tanques grandes y aislados, ignorando los clústeres densos.

\begin{figure}[h!]
\centering
\includegraphics[width=0.75\textwidth]{compare_019822ed-bf5b-4cc7-b911-4ee529d8f28b.jpg}
\caption{Comparativa cualitativa sobre una imagen de validación: anotación real
(verde), Hough tileado (azul), YOLOv8 (rojo).}
\label{fig:comparativa}
\end{figure}

\subsection{Líneas de mejora futuras}

\begin{itemize}
    \item \textbf{Señal de color real} en vez de solo brillo en escala de grises:
    los métodos deterministas actuales solo distinguen ``claro vs.\ oscuro'';
    incorporar el matiz/saturación (espacio HSV) podría separar tejados de tanque
    de otras superficies brillantes (hormigón, tejados residenciales) que hoy
    generan falsos positivos.
    \item \textbf{Cajas circulares explícitas} en el detector de red neuronal,
    siguiendo la propuesta de Wang et al.\ \cite{wang2025circular}, en vez de la
    caja rectangular genérica de YOLO --- debería mejorar la precisión de
    localización en tanques muy solapados.
    \item \textbf{Ensemble de métodos deterministas} (unión con NMS de \emph{blob
    detector} + \emph{watershed} universal): al usar lógicas de segmentación
    distintas sobre la misma señal de brillo, podrían tener falsos positivos no
    correlacionados y aciertos parcialmente complementarios.
    \item \textbf{Ajuste fino del lado YOLO}, que quedó deliberadamente fuera de
    alcance de este proyecto (se usó solo como referencia, no como objeto de
    optimización): barrido de umbral de confianza/NMS en inferencia, modelos más
    grandes (YOLOv8s/m), más épocas con aumentado de datos, o \emph{tiles} más
    pequeños específicamente para mejorar el \emph{recall} en los clústeres más
    densos.
    \item \textbf{Ampliación del dataset de entrenamiento}: con solo 78 imágenes
    de entrenamiento, un esquema de aprendizaje activo o la incorporación de más
    imágenes anotadas de la colección completa de Airbus (este es solo un
    subconjunto de demostración) sería la palanca más directa para seguir
    mejorando el modelo de referencia.
\end{itemize}

% =============================================================================
\section*{Anexo: reproducibilidad}
% =============================================================================
\addcontentsline{toc}{section}{Anexo: reproducibilidad}

Todo el código fuente, comentado y organizado en módulos (\texttt{src/}) y
\emph{scripts} ejecutables (\texttt{scripts/}), se entrega junto a esta memoria. La
bitácora experimental completa --- cada configuración probada, su resultado y el
razonamiento detrás de cada decisión --- está en el fichero \texttt{EXPERIMENTS.md}.
Las métricas crudas de cada experimento están en \texttt{results/*.json} y las
visualizaciones comparativas (anotación real vs.\ predicción) en
\texttt{results/figures/}.

% =============================================================================
% Bibliografía
% =============================================================================
\begin{thebibliography}{99}

\bibitem{duda1972hough}
Duda, R.\,O. \& Hart, P.\,E. (1972).
\emph{Use of the Hough Transformation to Detect Lines and Curves in Pictures}.
Communications of the ACM. (Fundamento teórico de \texttt{cv2.HoughCircles}).

\bibitem{faudi_yolox}
Faudi, J. --- \emph{Oil Storage Detection on Airbus Imagery with YOLOX}.
Medium / Artificialis.
\url{https://medium.com/artificialis/oil-storage-detection-on-airbus-imagery-with-yolox-9e38eb6f7e62}
Notebook asociado:
\url{https://www.kaggle.com/code/jeffaudi/oil-storage-detection-on-airbus-imagery-with-yolox}

\bibitem{ouzounis_maskrcnn}
Ouzounis, G. --- \emph{Oil-Storage Tank Instance Segmentation with Mask R-CNN}.
Medium.
\url{https://medium.com/@georgios.ouzounis/oil-storage-tank-instance-segmentation-with-mask-r-cnn-77c94433045f}

\bibitem{wang2025circular}
Wang et al. (2025).
\emph{Storage tank detection in remote sensing images based on circular bounding
boxes and large selective kernel}. Scientific Reports.
\url{https://www.nature.com/articles/s41598-025-27919-5}

\bibitem{beucher1993watershed}
Beucher, S. \& Meyer, F. (1993).
\emph{The morphological approach to segmentation: the watershed transformation}.
En \emph{Mathematical Morphology in Image Processing}, Marcel Dekker.
(Fundamento teórico de \texttt{cv2.watershed}).

\bibitem{lindeberg1998scale}
Lindeberg, T. (1998).
\emph{Feature Detection with Automatic Scale Selection}.
International Journal of Computer Vision.
(Fundamento teórico de la detección multiescala Laplaciano-de-Gaussiana/
Diferencia-de-Gaussianas implementada en \texttt{skimage.feature.blob\_dog}).

\bibitem{airbus_dataset}
Airbus DS Intelligence / Scale AI --- \emph{Airbus Oil Storage Detection Dataset}.
Kaggle. \url{https://www.kaggle.com/datasets/airbusgeo/airbus-oil-storage-detection-dataset}
(Licencia CC BY-NC-SA 4.0).

\bibitem{yolov8}
Jocher, G., Chaurasia, A. \& Qiu, J. (2023).
\emph{Ultralytics YOLOv8} [software].
\url{https://github.com/ultralytics/ultralytics}

\bibitem{bradski2000opencv}
Bradski, G. (2000).
\emph{The OpenCV Library}. Dr.\ Dobb's Journal of Software Tools.
(\texttt{cv2.SimpleBlobDetector}, transformadas y operaciones morfológicas
utilizadas en todo el proyecto).

\bibitem{vanderwalt2014scikit}
Van der Walt, S. et al. (2014).
\emph{scikit-image: image processing in Python}. PeerJ.
(\texttt{skimage.feature.peak\_local\_max}, \texttt{skimage.measure.regionprops},
\texttt{skimage.feature.blob\_dog}).

\end{thebibliography}

\end{document}
